\section{Conclusion}

\subsection{Summary}

This study developed a hybrid scheme in which a MEWMA chart accumulates a CAE reconstruction error and an OC-SVM score of the CAE features. In the primary experiment, on a simulated nonlinear process with five independent
datasets, every chart was calibrated empirically to an IC \ARL{} of 370. The learned components
behaved largely as intended, and on independent IC data the hybrid chart was somewhat more liberal,
and the classical benchmark somewhat more conservative, than the target. The hybrid chart had a lower $\ARL_1$ than both single-component ablations at every shift (within sampling error only against
$S(0.05)$ at $\delta\le0.4$), so the combined vector carries more usable information than either
learned signal alone, but it did not outperform the classical MEWMA: the two were
indistinguishable at $\delta=0.2$, and the benchmark was faster at every larger shift, with the hybrid chart needing 1.4 to 2.5 times as many post-change observations. Larger smoothing constants
shortened the hybrid chart's run-lengths modestly without closing the gap. In the external Extended TEP benchmark, with limits set by a false-alarm probability rather than an \ARL{}, both charts met their IC target and the classical MEWMA detected IDV(2) sooner at every severity, by 4.8 to 10.1 decisions on average.

\subsection{Recommendations}

First, for sustained mean shifts in a moderately nonlinear process and persistent step disturbances
such as TEP IDV(2), a well-calibrated classical MEWMA on the raw variables should remain the default, and a learned hybrid scheme should be preferred only where it demonstrably improves detection of the relevant disturbance structure under a common IC criterion; because both studies are simulations, this implication has yet to be confirmed on measured plant data. Second,
learned reconstruction and one-class information should be monitored jointly, since $H(0.05)$ had lower mean $\ARL_1$ than either ablation at every evaluated shift. Third, charts on overlapping windows should be
calibrated empirically on the dependent process, with delays reported both in monitoring
opportunities and in post-change observations.

\subsection{Future work}

The hybrid framework should be evaluated on faults that alter the shape, variability or dependence of the data
(variance changes, sensor faults, drifts and transients), including the other TEP faults, for which distance-type features may be better suited, and on other processes. Direction-preserving learned features, dependence-aware finite-start scaling for
overlapping windows, windowing that shortens the post-change ramp-in, and validation on measured
plant data also deserve study.
